\subsection{Modifying Conductance and Storage}

The internal flow properties---intercell conductance in the Node Property Flow
Package and storage in the Storage Package---are the simplest variables to modify
through the API, because MODFLOW 6 does not re-read them each stress period.
\texttt{NPF/CONDSAT} is computed once during Allocate and Read and then used
directly; \texttt{STO/SS} and \texttt{STO/SY} are read directly at every formulate
step. In all three cases a value set once persists for the rest of the
simulation, so it can be assigned a single time before the time loop and the
high-level \texttt{update} routine used to advance the simulation. Set
\texttt{CONDSAT} (not the hydraulic conductivity, which is inert after Allocate
and Read) to change conductance, and set \texttt{SS} or \texttt{SY} directly to
change storage.

\begin{lstlisting}[style=pseudocode]
initialize()
set MODEL/NPF/CONDSAT = ...       # change intercell conductance (not K, which is inert)
set MODEL/STO/SS = ...            # change confined (specific) storage
set MODEL/STO/SY = ...            # change water-table storage (specific yield)
while time < end_time:
    update()                     # nothing re-reads these, so the change persists
finalize()
\end{lstlisting}

\noindent A property can also be changed partway through a simulation by setting
it before the time step from which the new value should take effect; it then
persists until it is changed again.

\begin{lstlisting}[style=pseudocode]
initialize()
while time < end_time:
    if time >= change_time:
        set MODEL/NPF/CONDSAT = ...   # takes effect from this step onward
    update()
finalize()
\end{lstlisting}

\noindent The set-once approach is valid only when MODFLOW 6 does not rewrite the
variable. If time-varying hydraulic conductivity (TVK), time-varying storage
(TVS), or the Viscosity (VSC) Package is active, the corresponding variable is
refreshed every time step and must instead be set after \texttt{prepare\_time\_step}
on each step, using the pattern shown next for stress-package input. Similarly, if
the Horizontal Flow Barrier (HFB) Package is active, it overwrites \texttt{CONDSAT}
on the connections it spans at every Read and Prepare, so those entries must be
re-applied each step rather than set once.

\subsection{Updating Stress-Package Input at Run Time}

The values in a stress package are written by MODFLOW 6 inside
\texttt{prepare\_time\_step}, which performs Read and Prepare (only on the first
time step of a stress period that reads a new period-data block) and Advance
(every time step). A value supplied as a constant in the package file is written
during Read and Prepare; a value supplied by a time series is written during
Advance. An API override must therefore be applied \emph{after}
\texttt{prepare\_time\_step}; a value set before it is overwritten by whichever
of those two steps writes the variable.

In the simplest case the value can be set before the time step and the step
advanced with the high-level convenience routine \texttt{update}, which runs an
entire time step (\texttt{prepare\_time\_step}, the iteration loop, and
\texttt{finalize\_time\_step}) in a single call. This works only when MODFLOW 6
will not rewrite the value during the step. When it would---most commonly for
time-series values---the lower-level routines are used instead so the value can
be set after \texttt{prepare\_time\_step}. The pseudo-code uses the memory
address \texttt{MODEL/PKG/COND} to stand for any stress-package input variable.

\subsubsection{Simplest use case: setting a value before the time step}

A value can be set before the time step (and the step advanced with
\texttt{update}) when MODFLOW 6 will not rewrite it during that step. This is the
case only when both of the following hold: the value is not supplied by a time
series (a time series rewrites it every step during Advance), and the step is not
the first step of a stress period that reads a new period-data block (Read and
Prepare rewrites the named array then). Mid-period time steps, and stress periods
that reuse the previous period's data, satisfy both conditions.

\begin{lstlisting}[style=pseudocode]
initialize()
while time < end_time:
    # valid only when, for this step, MF6 will not overwrite the value:
    #   - the value is not driven by a time series, AND
    #   - this is not the first step of a stress period with a new PERIOD block
    if value_will_not_be_overwritten_this_step:
        set MODEL/PKG/COND = ...
    update()                     # = prepare_time_step + iteration loop + finalize_time_step
finalize()
\end{lstlisting}

\subsubsection{Standard pattern: update after \texttt{prepare\_time\_step}}

When the value may be rewritten during the step---most importantly a time-series
value, which is refreshed every step during Advance---it must be set after
\texttt{prepare\_time\_step}. Setting it every time step at that point, using the
lower-level routines, is the safe and general pattern for both package-file and
time-series values.

\begin{lstlisting}[style=pseudocode]
initialize()
while time < end_time:
    prepare_time_step(dt)        # Read and Prepare (period start) + Advance (time series)
    set MODEL/PKG/COND = ...      # safe here: after MF6 has written its own value
    prepare_solve()
    repeat:
        converged = solve()      # one outer (Picard/Newton) iteration
    until converged
    finalize_solve()
    finalize_time_step()
finalize()
\end{lstlisting}

\subsubsection{Head-dependent values within the iteration loop}

Each \texttt{solve} call performs one outer (Picard or Newton) iteration, in which
the stress packages are formulated from their input variables. Modifying a
package input before each \texttt{solve} therefore re-linearizes the boundary at
every iteration, which can be used to implement a value that depends on the
current head (a custom head-dependent boundary or nonlinear coefficient). Modify
the package input (for example \texttt{COND}); MODFLOW 6 recomputes the
corresponding \texttt{HCOF} and \texttt{RHS} terms from it on each iteration. The
outer iteration loop converges when the head and the value it depends on reach a
consistent state.

\begin{lstlisting}[style=pseudocode]
initialize()
while time < end_time:
    prepare_time_step(dt)
    prepare_solve()
    repeat:                              # outer (Picard/Newton) iteration loop
        head = get_value_ptr("MODEL/X")
        set MODEL/PKG/COND = f(head)     # modify the package input before each solve
        converged = solve()             # this iteration formulates with the new value
    until converged                     # head and value reach self-consistency
    finalize_solve()
    finalize_time_step()
finalize()
\end{lstlisting}

A complete, runnable version of the standard pattern---with one general-head
boundary whose conductance is a constant in the package file and a second whose
conductance is supplied by a time series---is provided in
\texttt{doc/mf6api/examples/update\_stress\_via\_api.py}.
